\ExamPart{Câu tự luận}{3}{ }

\NQuestion Giải bất phương trình $\left(x^2-4x+3\right)\left(x^2-1\right)\le 0$.
\AnswerLine{$[-1;3]$}
\SolutionBlock{Ta có
\[
\left(x^2-4x+3\right)\left(x^2-1\right)=\left(x-1\right)\left(x-3\right)\left(x-1\right)\left(x+1\right)=\left(x-1\right)^2\left(x-3\right)\left(x+1\right).
\]
Vì $\left(x-1\right)^2\ge0$ với mọi $x$, nên dấu của biểu thức phụ thuộc vào $\left(x-3\right)\left(x+1\right)$, đồng thời các điểm làm biểu thức bằng $0$ đều được nhận do dấu ``$\le0$''. Ta giải:
\[
\left(x-3\right)\left(x+1\right)\le0 \Leftrightarrow -1\le x\le3.
\]
Do đó tập nghiệm là $[-1;3]$.}

\NQuestion Tìm tất cả các giá trị của $m$ để phương trình $x^2-2\left(m+1\right)x+m^2-2m-3=0$ có hai nghiệm phân biệt $x_1,x_2$ thỏa mãn $x_1+x_2>2$ và $x_1x_2<0$.
\AnswerLine{$0<m<3$}
\SolutionBlock{Theo hệ thức Viète:
\[
x_1+x_2=2\left(m+1\right),\qquad x_1x_2=m^2-2m-3.
\]
Điều kiện $x_1+x_2>2$ cho
\[
2\left(m+1\right)>2 \Leftrightarrow m>0.
\]
Điều kiện $x_1x_2<0$ cho
\[
m^2-2m-3<0 \Leftrightarrow \left(m-3\right)\left(m+1\right)<0 \Leftrightarrow -1<m<3.
\]
Nếu $x_1x_2<0$ thì phương trình chắc chắn có hai nghiệm thực phân biệt. Kết hợp với $m>0$, ta được
\[
0<m<3.
\]}

\NQuestion Trong mặt phẳng tọa độ $Oxy$, cho ba điểm $A\left(1;2\right),\ B\left(5;4\right),\ C\left(3;-2\right)$.
\begin{SubQuestions}
  \item Tính tọa độ các vectơ $\overrightarrow{AB}$ và $\overrightarrow{AC}$.
  \item Chứng minh tam giác $ABC$ vuông tại $A$.
  \item Viết phương trình đường thẳng $BC$.
  \item Tính khoảng cách từ điểm $A$ đến đường thẳng $BC$.
\end{SubQuestions}
\AnswerLine{$\overrightarrow{AB}=\left(4;2\right),\ \overrightarrow{AC}=\left(2;-4\right),\ BC:3x-y-11=0,\ d\left(A,BC\right)=\sqrt{10}$}
\SolutionBlock{\begin{SubQuestions}
\item Ta có
\[
\overrightarrow{AB}=\left(5-1;4-2\right)=\left(4;2\right),\qquad \overrightarrow{AC}=\left(3-1;-2-2\right)=\left(2;-4\right).
\]
\item Ta tính tích vô hướng:
\[
\overrightarrow{AB}\cdot\overrightarrow{AC}=4\cdot2+2\cdot\left(-4\right)=8-8=0.
\]
Suy ra $\overrightarrow{AB}\perp\overrightarrow{AC}$, nên tam giác $ABC$ vuông tại $A$.
\item Đường thẳng $BC$ đi qua $B\left(5;4\right)$ và $C\left(3;-2\right)$ có hệ số góc
\[
m=\dfrac{-2-4}{3-5}=3.
\]
Phương trình qua $B$:
\[
y-4=3\left(x-5\right)\Leftrightarrow 3x-y-11=0.
\]
\item Khoảng cách từ $A\left(1;2\right)$ đến $BC:3x-y-11=0$ là
\[
d\left(A,BC\right)=\dfrac{|3\cdot1-2-11|}{\sqrt{3^2+\left(-1\right)^2}}=\dfrac{10}{\sqrt{10}}=\sqrt{10}.
\]
\end{SubQuestions}}

\NQuestion Trong mặt phẳng tọa độ $Oxy$, cho điểm $M\left(2;-1\right)$ và đường thẳng $d:2x-y+3=0$.
\begin{SubQuestions}
  \item Viết phương trình đường thẳng $\Delta$ đi qua $M$ và song song với $d$.
  \item Viết phương trình đường thẳng $\Delta'$ đi qua $M$ và vuông góc với $d$.
  \item Tìm tọa độ hình chiếu vuông góc $H$ của $M$ trên $d$.
\end{SubQuestions}
\AnswerLine{$\Delta:2x-y-5=0,\ \Delta':x+2y=0,\ H\left(-\dfrac65;\dfrac35\right)$}
\SolutionBlock{\begin{SubQuestions}
\item Vì $\Delta\parallel d$ nên $\Delta$ có dạng
\[
2x-y+c=0.
\]
Thay $M\left(2;-1\right)$ vào:
\[
2\cdot2-\left(-1\right)+c=0 \Rightarrow c=-5.
\]
Vậy $\Delta:2x-y-5=0$.
\item Đường thẳng $d$ có hệ số góc $2$, nên đường thẳng vuông góc với $d$ có hệ số góc $-\dfrac12$. Phương trình qua $M\left(2;-1\right)$ là
\[
y+1=-\dfrac12\left(x-2\right)\Leftrightarrow x+2y=0.
\]
\item Tọa độ $H$ là giao điểm của
\[
d:2x-y+3=0,\qquad \Delta':x+2y=0.
\]
Từ $\Delta'$ suy ra $x=-2y$. Thay vào $d$:
\[
2\left(-2y\right)-y+3=0 \Rightarrow -5y+3=0 \Rightarrow y=\dfrac35.
\]
Khi đó
\[
x=-2\cdot\dfrac35=-\dfrac65.
\]
Vậy $H\left(-\dfrac65;\dfrac35\right)$.
\end{SubQuestions}}
